\pdfminorversion=4
%% Safe RTO under unmeasured disturbances -- IEEE TCST submission (double-column, <=12 pp)
%% Compile on Overleaf: pdflatex -> bibtex -> pdflatex x2. IEEEtran.cls + IEEEtran.bst are native there.
\documentclass[journal]{IEEEtran}

\usepackage[T1]{fontenc}
\usepackage{amsmath,amssymb,amsthm}
\usepackage{graphicx}
\usepackage{booktabs}
\usepackage{url}
\usepackage[numbers,sort&compress]{natbib}
\graphicspath{{figures/}}

\newtheorem{theorem}{Theorem}
\newtheorem{corollary}{Corollary}
\newtheorem{definition}{Definition}
\newtheorem{assumption}{Assumption}

\begin{document}

\title{Safe Real-Time Optimization Under Unmeasured Disturbances:\\
A Finite-Sample, Distribution-Free Constraint-Satisfaction Guarantee}

\author{Bien~Don~Busico%
\thanks{B. D. Busico is with Mapua Malayan College Mindanao, Davao City, Philippines
(e-mail: bienbusico@gmail.com).}}

\markboth{IEEE Transactions on Control Systems Technology}%
{Busico: Safe Real-Time Optimization Under Unmeasured Disturbances}

\maketitle

\begin{abstract}
Real-time optimization (RTO) drives a process against a product-quality constraint it cannot
measure online, so it must subtract a back-off from the specification. Conformal prediction is
an attractive, distribution-free way to size that back-off from data, yet we prove that the
guarantee it provides does not survive the optimization. We show that an RTO respecting a
marginally valid conformal back-off carries no coverage guarantee at the operating point it
selects: the realized constraint-violation probability can exceed the nominal level by an
arbitrary factor, because the optimizer settles on the active boundary where the margin is
least conservative relative to the local risk. This is a constraint-space analogue of the
winner's curse. We then prove that conditional calibration restores a finite-sample,
distribution-free guarantee. An a-posteriori recalibration at the selected point achieves exact
validity, and, when fresh calibration data at that point are unavailable, a calibration made
uniform over the decision set guarantees a realized violation no worse than the nominal level
plus a term of order $\sqrt{d/n}$ that is small for the low-dimensional decisions typical of
RTO. On a rigorous DWSIM Peng-Robinson debutanizer twin, a marginally valid back-off realizes
roughly five times the nominal violation, the conditional method tracks a plant-aware oracle at
near-oracle profit, and the a-posteriori inflation factor is an interpretable data-adequacy
diagnostic that bounds the safe operating regime.
\end{abstract}

\begin{IEEEkeywords}
Real-time optimization, chance-constrained control, conformal prediction, safe learning,
constraint back-off, distribution-free guarantees, process control.
\end{IEEEkeywords}

\IEEEpeerreviewmaketitle

\section{Introduction}
\label{sec:intro}

\IEEEPARstart{R}{eal-time} optimization (RTO) recomputes the operating point of a process to
maximize an economic objective as conditions drift. Its value comes from pushing the process
against its constraints; a distillation column earns more as it drives a product composition to
the edge of specification. The constraints that matter most, however, are often on quality
variables that are not measured online. A laboratory assay or an inferential estimate stands in
for a true composition the controller never sees in real time. To operate safely against such a
constraint, an RTO must subtract a \emph{back-off} from the specification: a margin that absorbs
the gap between the model's nominal prediction and the realized quality under an unmeasured
disturbance. Sizing that margin is the central difficulty of RTO under uncertainty. Set it too
small and the process violates specification whenever the disturbance is adverse; set it too
large and the economic benefit of the optimization is given away.

Conformal prediction is an appealing tool for the job. It turns any predictor into an interval
with a finite-sample, distribution-free coverage guarantee, assuming only that the data are
exchangeable, with no Gaussian residuals and no parametric disturbance law
\cite{vovk2005,lei2018,angelopoulos2023}. Taking the upper edge of a conformal interval as the
back-off appears to deliver exactly what a chance constraint asks: coverage of the realized
quality at the nominal level, learned from data. This paper shows that the guarantee does not
survive contact with the optimizer, and it supplies a back-off construction, with proof, that
restores it.

The failure is structural. A conformal margin guarantees coverage \emph{on average} over the
calibration distribution of operating points, but the optimizer does not sample that
distribution. It deliberately seeks the boundary of the feasible region, where the margin is
tightest relative to the local risk and therefore where coverage is worst. The argmax is a
selection operator, and it selects for under-coverage. The conformal-predictive-programming
literature \cite{zhao2024cpp} observed abstractly that a decision optimized against sampled
constraints is no longer independent of them; our concern is the concrete consequence for a
control layer, namely a safety hazard of provable and quantified severity, and the construction
that removes it. Theorem~\ref{thm:vacuous} shows that marginal validity implies \emph{no}
nontrivial bound on the realized violation at the selected operating point. This is a
constraint-space analogue of the winner's curse in post-selection inference.

What repairs it is conditional validity, and we prove two guarantees for it
(Theorem~\ref{thm:repair}). A back-off recalibrated a-posteriori at the selected operating point,
with an independent calibration batch, achieves exact finite-sample validity because the
selection is then independent of the calibration data. When such a batch is unavailable and the
same historical data must serve both optimization and calibration, a calibration made uniform
over the decision set restores validity up to a term of order $\sqrt{d/n}$, which is small for
the low-dimensional decisions characteristic of RTO. A back-off built from conformalized quantile
regression \cite{romano2019} is the practical route to this conditional target, and the
a-posteriori inflation factor it requires is an observable diagnostic of calibration adequacy
(Corollary~\ref{cor:kappa}).

We demonstrate the theory on a rigorous DWSIM Peng-Robinson debutanizer twin. A marginally valid
back-off, including a locally adaptive heteroscedastic one, realizes roughly five times the
nominal violation across the disturbance range; a plant-aware oracle holds the target by
construction, confirming that the failure lies in the back-off and not in the chance-constrained
formulation; the conditional method tracks the oracle at near-oracle profit; and the inflation
factor climbs from near unity to nearly sixfold before the method exhausts its calibration data,
bounding the safe operating regime. We are explicit that the contribution is calibrated safety,
not profit: at a well-controlled feed every method earns within half a percent of the
deterministic optimum, so the value of the method is the violation it removes.

The contributions are as follows.
\begin{itemize}
\item \textbf{C1.} We prove that optimizing against a marginally valid conformal back-off carries
no coverage guarantee at the selected operating point: the realized violation can exceed the
nominal level by an arbitrary factor (Theorem~\ref{thm:vacuous}).
\item \textbf{C2.} We prove that conditional calibration restores a finite-sample,
distribution-free guarantee: a-posteriori recalibration achieves exact validity, and uniform
calibration guarantees a realized violation of at most the nominal level plus an
$O(\sqrt{d/n})$ term (Theorem~\ref{thm:repair}), and we tie the required inflation factor to
calibration adequacy (Corollary~\ref{cor:kappa}).
\item \textbf{C3.} We validate the theory on a rigorous distillation RTO twin, quantifying the
fivefold failure of marginal back-offs, the near-oracle safety and profit of the conditional
method, and the inflation factor as a data-adequacy diagnostic that bounds the safe regime.
\end{itemize}

The remainder of the paper is organized as follows. Section~\ref{sec:methods} formalizes the
chance-constrained RTO and the conformal back-offs. Section~\ref{sec:theory} proves the selection
effect and its repair. Section~\ref{sec:case} describes the twin and the experimental design,
Section~\ref{sec:results} reports the results, and Section~\ref{sec:discussion} draws the design
rule, the limitations, and the bridge to a closed-loop realization.

\subsection{Related work}
\label{sec:related}
RTO under uncertainty classically addresses plant-model mismatch through modifier adaptation,
which corrects the optimizer with measured gradients \cite{marchetti2009,chachuat2009}, and
through chance-constrained programming, which bounds constraint violation under an assumed
disturbance law \cite{li2008,bradford2019,bradford2020}. Both presuppose a model of the
uncertainty. Conformal prediction removes that assumption, guaranteeing finite-sample coverage
under exchangeability alone \cite{vovk2005,lei2018,angelopoulos2023}, with conditional and
covariate-shift extensions \cite{romano2019,gibbs2021,gibbscherian2023,tibshirani2019,barber2023}.
Its use inside optimization is recent \cite{zhao2024cpp,zhang2024}: conformal predictive
programming recognizes that optimization couples the decision to the sampled constraints. We
contribute the control-theoretic consequence of that coupling, a proof that marginal back-offs
are unsafe at the optimum, and a proven repair with an accompanying diagnostic.

\section{Problem Formulation and Conformal Back-Offs}
\label{sec:methods}

\subsection{Chance-constrained RTO under an unmeasured disturbance}
\label{sec:methods-rto}
We consider steady-state RTO of a distillation column whose product-quality constraint cannot be
measured online (Fig.~\ref{fig:rto}). The decision is the operating point $u=(R,D)$, with $R$ the
reflux ratio and $D$ the distillate molar flow; the admissible set is the box
$\mathcal U=[0.8,3.0]\times[33,37]$, with $D$ in kmol\,h$^{-1}$. The plant is subject to an
exogenous disturbance $z$, the feed light-key (n-butane) mole fraction, neither manipulated nor
measured in real time, modeled as a truncated normal
$z\sim\mathcal F_\sigma=\mathcal N(z_0,\sigma_z^2)|_{[z_{\mathrm{lo}},z_{\mathrm{hi}}]}$ with
$z_0=0.35$, $[z_{\mathrm{lo}},z_{\mathrm{hi}}]=[0.30,0.40]$, and $\sigma_z$ the disturbance
magnitude (swept in Section~\ref{sec:case}; the operationally realistic value is
$\sigma_z\approx0.006$).

The constrained quantity is the bottoms light-key composition $g(u,z)=x_B(R,D,z)$, which must
satisfy $g(u,z)\le\bar g$ with $\bar g=0.02$. Because $g$ is unmeasured, the optimizer sees only a
nominal model $\hat g(u)=x_B(R,D,z_0)$. The economic objective $\pi(u)$ is an operating profit
with two product streams valued by mass less the reboiler energy
(Section~\ref{sec:methods-twin}). The principled formulation is a chance constraint,
\begin{equation}
\max_{u\in\mathcal U}\;\pi(u)\quad\text{s.t.}\quad
\mathbb{P}_{z\sim\mathcal F_\sigma}\!\left[\,g(u,z)\le\bar g\,\right]\ge 1-\alpha,
\label{eq:chance}
\end{equation}
with miss level $\alpha=0.10$. Writing $Q^z_{1-\alpha}[g(u,\cdot)]$ for the $(1-\alpha)$-quantile
of $g(u,z)$, \eqref{eq:chance} is equivalent to $Q^z_{1-\alpha}[g(u,\cdot)]\le\bar g$. RTO has no
access to this quantile, so it enforces a surrogate by subtracting a back-off $C(u)\ge0$ from the
spec applied to the nominal model,
\begin{equation}
\max_{u\in\mathcal U}\;\pi(u)\quad\text{s.t.}\quad \hat g(u)+C(u)\le\bar g .
\label{eq:backoff}
\end{equation}
The entire content of an uncertainty-aware RTO lies in how $C(u)$ is constructed: the back-off is
correct when $\hat g(u)+C(u)\ge Q^z_{1-\alpha}[g(u,\cdot)]$, so that any $u$ feasible in
\eqref{eq:backoff} is feasible in \eqref{eq:chance}.

A benchmark is the oracle. Because $g(u,z)$ is monotone increasing in $z$ at every $u$ (verified
on the twin), the conditional quantile collapses to $Q^z_{1-\alpha}[g(u,\cdot)]=g(u,z_{1-\alpha})$
with $z_{1-\alpha}=\mathcal F_\sigma^{-1}(1-\alpha)$, giving the oracle back-off
$C_{\mathrm{orc}}(u)=[g(u,z_{1-\alpha})-\hat g(u)]_+$. The oracle needs the true plant response and
is not deployable; it makes the realized violation at any active optimum equal to $\alpha$, which
we use to confirm the chance-constraint machinery is sound.

\subsection{Conformal constraint back-offs}
\label{sec:methods-backoffs}
Assume a calibration set of plant evaluations $\{(u_i,z_i,g_i)\}_{i=1}^n$ with
$u_i\sim\mathrm{Unif}(\mathcal U)$, $z_i\sim\mathcal F_\sigma$, $g_i=g(u_i,z_i)$; in deployment
these are historical operating records. Define residuals $r_i=g_i-\hat g(u_i)$. For a one-sided
upper bound at level $1-\alpha$, the split-conformal quantile of scores $\{s_i\}$ is the order
statistic
\begin{equation}
\widehat Q_{1-\alpha}(\{s_i\})=s_{(k)},\qquad k=\lceil(n+1)(1-\alpha)\rceil,
\label{eq:quantile}
\end{equation}
which is $+\infty$ if $k>n$. We compare four constructions of $C(u)$, all floored at zero. The
\emph{fixed} margin ignores the operating point, $C_{\mathrm{fix}}=\widehat Q_{1-\alpha}(\{r_i\})$;
it is marginally valid but spends the same margin everywhere. The \emph{normalized}, locally
adaptive back-off scales the score by a Gaussian-process spread estimate $\hat s(u)$ fit to
$\{|r_i|\}$ \cite{lei2018},
$C_{\mathrm{nrm}}(u)=\widehat Q_{1-\alpha}(\{r_i/\hat s(u_i)\})\,\hat s(u)$; it is the back-off a
soft-sensor interval supplies, marginally but not conditionally valid. \emph{Conformalized
quantile regression} (CQR) \cite{romano2019} is the conditional construction: fit a
gradient-boosted quantile regressor $\hat q_{1-\alpha}(u)$, form scores
$E_i=g_i-\hat q_{1-\alpha}(u_i)$, and set
$C_{\mathrm{cqr}}(u)=[\hat q_{1-\alpha}(u)+\widehat Q_{1-\alpha}(\{E_i\})-\hat g(u)]_+$. Because
$\hat q_{1-\alpha}(u)$ estimates the conditional quantile at each $u$, $C_{\mathrm{cqr}}$ tracks
the local risk, and the conformal correction restores finite-sample marginal validity on top.

\section{The Selection Effect and Its Repair}
\label{sec:theory}

This section makes precise the failure of marginal back-offs under optimization and proves that
conditional calibration repairs it. Throughout, let
$M(u)\triangleq\mathbb{P}_{z\sim\mathcal F_\sigma}[g(u,z)>\hat g(u)+C(u)]$ be the realized
\emph{miss} of the back-off at $u$. Because any $u$ feasible in \eqref{eq:backoff} satisfies
$\hat g(u)+C(u)\le\bar g$, the event $\{g(u,z)>\bar g\}$ is contained in
$\{g(u,z)>\hat g(u)+C(u)\}$, so the specification-violation probability is bounded by the miss,
$\mathbb P_z[g(u,z)>\bar g]\le M(u)$. Controlling $M$ at the selected point controls the chance
constraint \eqref{eq:chance}.

\begin{definition}[Marginal and conditional validity]
A back-off $C$ is \emph{marginally valid} at level $\alpha$ over a reference distribution $\rho$
on $\mathcal U$ if $\mathbb E_{u\sim\rho}[M(u)]\le\alpha$. It is \emph{conditionally valid} at
level $\alpha$ on $\mathcal S\subseteq\mathcal U$ if $M(u)\le\alpha$ for every $u\in\mathcal S$.
\end{definition}

Split and normalized conformal back-offs are marginally valid by construction
\cite{vovk2005,lei2018}, but the conditional miss varies over $\mathcal U$
\cite{gibbscherian2023}. The next result shows this gap is a safety failure once an optimizer is
in the loop.

\begin{theorem}[Marginal validity is vacuous at the optimum]
\label{thm:vacuous}
Fix $\alpha\in(0,1)$. For every $\kappa\in[1,1/\alpha]$ there is an instance $(g,\hat g,\pi,C)$
with $C$ marginally valid at level $\alpha$ over $\rho=\mathrm{Unif}(\mathcal U)$ for which the RTO
solution $u^\star=\arg\max_{u\in\mathcal U}\{\pi(u):\hat g(u)+C(u)\le\bar g\}$ satisfies
$M(u^\star)=\kappa\alpha$. Consequently marginal validity alone implies no bound on $M(u^\star)$
tighter than the trivial $M(u^\star)\le1$.
\end{theorem}

\begin{proof}
Partition $\mathcal U=A\cup B$ with $\rho(A)=1/\kappa\in[\alpha,1]$, admissible since
$\kappa\in[1,1/\alpha]$. The miss $M(u)$ is continuous and strictly decreasing in $C(u)$ from
$\mathbb P_z[g(u,z)>\hat g(u)]$ to $0$; choose $C$ so that $M(u)=\kappa\alpha$ on $A$ (attainable
whenever $\kappa\alpha\le1$) and $M(u)=0$ on $B$. Then
$\mathbb E_{u\sim\rho}[M(u)]=\rho(A)\kappa\alpha+\rho(B)\cdot0=\alpha$, so $C$ is marginally valid.
Place the profit so its constrained maximizer lies in $A$: set $\hat g(u)+C(u)=\bar g-\epsilon$ on
$A$ with $\pi$ increasing toward that boundary and strictly smaller on $B$. Then $u^\star\in A$ and
$M(u^\star)=\kappa\alpha$. As $\kappa$ ranges over $[1,1/\alpha]$ with $\kappa\alpha\le1$, no bound
better than $M(u^\star)\le1$ follows.
\end{proof}

Theorem~\ref{thm:vacuous} is the winner's curse in constraint space. The optimizer takes, among
points feasible under the back-off, the most profitable, which presses hardest on the spec where
the margin is tightest relative to the local risk; averaging coverage over $\mathcal U$ hides that
point. The fivefold inflation reported in Section~\ref{sec:results} is one realization, and it is
not cured by an adaptive width, because a smooth heteroscedastic scale still leaves pockets of
under-coverage the optimizer finds. The repair is to control the conditional miss at the selected
point. We give two guarantees, according to whether fresh calibration data can be obtained there.

\begin{assumption}[Exchangeable calibration]
\label{ass:exch}
For a fixed operating point $u$, the calibration disturbances and the deployment disturbance are
i.i.d.\ draws from $\mathcal F_\sigma$. Residuals are almost surely finite.
\end{assumption}

\begin{theorem}[Conditional calibration restores finite-sample validity]
\label{thm:repair}
Let $u^\star$ be produced by any selection rule using a dataset $D_n$.
\begin{enumerate}
\item[(a)] \emph{A-posteriori recalibration.} Draw a fresh sample
$z_1',\dots,z_m'\stackrel{\text{iid}}{\sim}\mathcal F_\sigma$ independent of $D_n$, and set the
back-off at $u^\star$ to the split-conformal quantile \eqref{eq:quantile} of the fresh residuals
$\{g(u^\star,z_j')-\hat g(u^\star)\}_{j=1}^m$. Then under Assumption~\ref{ass:exch},
\[
\mathbb P\big[g(u^\star,z)\le\hat g(u^\star)+C(u^\star)\big]\ge
\frac{\lceil(1-\alpha)(m+1)\rceil}{m+1}\ge 1-\alpha,
\]
so $M(u^\star)\le\alpha$, distribution-free and finite-sample.
\item[(b)] \emph{Data-reuse, uniform calibration.} Suppose instead the same $D_n$ builds
$C(\cdot)$ within a back-off class $\mathcal C$ of covering number $N(\mathcal C,\varepsilon)$ and
selects $u^\star$. Calibrate $C$ so the empirical miss on $D_n$ is at most $\alpha-\Delta$
uniformly over $\mathcal U$, with
\[
\Delta=\sqrt{\frac{\log N(\mathcal C,\varepsilon)+\log(2/\delta)}{2n}}+\varepsilon L,
\]
where $L$ is a Lipschitz constant of $u\mapsto M(u)$. Then with probability at least $1-\delta$
over $D_n$, $M(u^\star)\le\alpha$ simultaneously for every $u^\star\in\mathcal U$, hence for the
data-selected optimum.
\end{enumerate}
\end{theorem}

\begin{proof}
(a) Conditional on $D_n$, the point $u^\star$ is fixed and independent of the fresh sample, so the
residuals $R_j'=g(u^\star,z_j')-\hat g(u^\star)$ and the deployment residual
$R'=g(u^\star,z)-\hat g(u^\star)$ are i.i.d.\ by Assumption~\ref{ass:exch}. Split-conformal
validity \cite{vovk2005,lei2018} gives $\mathbb P[R'\le R_{(k)}']\ge k/(m+1)$ with
$k=\lceil(1-\alpha)(m+1)\rceil$; taking $C(u^\star)=R_{(k)}'$ and complementing gives the claim.
The selection is independent of the calibration sample and does not enter.

(b) Let $\widehat M(u)=\frac1n\sum_{i=1}^n\mathbf 1\{g(u_i,z_i)>\hat g(u_i)+C(u_i)\}$ be the
empirical miss and $M(u)$ its population value. Over an $\varepsilon$-cover of $\mathcal C$ of size
$N(\mathcal C,\varepsilon)$, Hoeffding's inequality and a union bound give, with probability
$\ge1-\delta$, $\sup_{C\in\mathcal C}|\widehat M-M|\le\sqrt{(\log N(\mathcal C,\varepsilon)+
\log(2/\delta))/(2n)}$; Lipschitz continuity of $M$ adds the discretization term $\varepsilon L$,
so $\sup_{u\in\mathcal U}|\widehat M(u)-M(u)|\le\Delta$. Calibrating $\widehat M(u)\le\alpha-\Delta$
for all $u$ yields $M(u)\le\alpha$ for all $u$, hence for $u^\star\in\mathcal U$ regardless of how
it was chosen from $D_n$.
\end{proof}

For the low-dimensional decisions typical of RTO, $u\in\mathbb R^d$ with small $d$ and a
finite-dimensional back-off class, $\log N(\mathcal C,\varepsilon)=O(d\log(1/\varepsilon))$, so
optimizing $\varepsilon$ gives $\Delta=\tilde O(\sqrt{d/n})$: the uniform guarantee costs only a
mild inflation of the calibration level (here $d=2$). Part~(a) is the tighter route when the plant
can be run at $u^\star$ to collect a calibration batch before committing; part~(b) is the
guarantee when only historical design data are available. The CQR back-off of
Section~\ref{sec:methods-backoffs} is the practical instantiation, targeting the conditional
quantile so the systematic selection gap of Theorem~\ref{thm:vacuous} vanishes and only the
finite-sample term remains.

\begin{corollary}[The inflation factor as a data-adequacy diagnostic]
\label{cor:kappa}
Let $C_0$ be the conditional back-off and
$\kappa^\star=\min\{\kappa\ge1:M(u^\star(\kappa))\le\alpha\}$ the smallest multiplier of $C_0$ that
certifies the target at the resolved optimum
$u^\star(\kappa)=\arg\max_u\{\pi(u):\hat g(u)+\kappa C_0(u)\le\bar g\}$. Then $\kappa^\star\to1$ when
$C_0$ is already conditionally valid (Theorem~\ref{thm:repair} with $\Delta\to0$), $\kappa^\star$
increases as the calibration sample shrinks relative to the disturbance spread, and the instance
is certifiably infeasible when no finite $\kappa$ attains $M\le\alpha$. Thus $\kappa^\star$ is an
observable surrogate for the calibration adequacy $\Delta$.
\end{corollary}

Operationally, $\kappa^\star$ is found by the a-posteriori search: after solving \eqref{eq:backoff}
with $C_0$, inflate the back-off by $\kappa$, re-solve, and take the smallest $\kappa$ whose
resolved optimum meets the target on held-out data, using a fixed validation sample so the search
is monotone and deterministic and an independent test sample for the reported violation; when even
the largest feasible $\kappa$ violates, the instance is declared infeasible, the honest outcome
when the data cannot certify the target. Section~\ref{sec:results} shows $\kappa^\star$ climbing
from near unity to nearly sixfold before the conditional method exhausts its calibration data,
which is the empirical face of Corollary~\ref{cor:kappa}.

\begin{figure*}[t]
\centering
\includegraphics[width=\textwidth]{F1_rto_loop.png}
\caption{Chance-constrained RTO with a conformal constraint back-off. The nominal model and the
back-off feed the optimizer, which sets $(R,D)$ on the twin. The constrained bottoms composition
$x_B$ is not measured online; a closed-loop soft-sensor feedback path (grey, dashed) is left to
future work.}
\label{fig:rto}
\end{figure*}

\section{Process Twin and Experimental Design}
\label{sec:case}

The plant is a rigorous DWSIM 9.0.5 model of a debutanizer \cite{dwsim}, reduced to the binary
n-butane/n-hexane system (light key and stabilized-gasoline proxy) under the Peng-Robinson
equation of state \cite{peng1976}, with nine equilibrium stages (condenser indexed $0$, reboiler
$8$), feed on stage $4$, top pressure $4.7$\,bar, a $0.4$\,bar column drop, and feed
$100$\,kmol\,h$^{-1}$. The model is driven from Python through DWSIM's automation interface; a
feed-composition campaign over a grid in $(z,R,D)$ yields the plant evaluations. The binary
reduction keeps every twin state checkable against first-principles vapor-liquid equilibrium and
matches the soft-sensor physics proxy of the companion work \cite{companion2026}.

Two Gaussian-process surfaces stand in for the plant in \eqref{eq:backoff}. The truth surface
$g(u,z)=x_B(R,D,z)$ is fit on the full campaign and supplies the realized bottoms composition used
to score violations and define the oracle; the nominal surfaces $\hat g(u)$, distillate purity,
and reboiler duty are fit on the nominal ($z_0=0.35$) sweep and define the optimizer's model and
objective. All GPs use a Mat\'ern-$5/2$ kernel on standardized inputs with a
maximum-a-posteriori regularization of the log-hyperparameters and a fixed optimizer seed for
reproducibility. Operating profit values the overhead at the light-key (LPG) price and the bottoms
at the stabilized-gasoline price, both by mass, less the reboiler energy cost; because gasoline
outvalues the light key per unit mass, profit rises as the column rides the bottoms-composition
spec from below, making the constraint economically active and the back-off width a direct cost.

We sweep the disturbance magnitude $\sigma_z\in\{0.006,0.010,0.015,0.020,0.025\}$, from the
operationally realistic $0.006$ to more than four times that, and for each solve \eqref{eq:backoff}
with the fixed, normalized, CQR, and oracle back-offs at $\alpha=0.10$. Realized violation is a
Monte-Carlo estimate on an independent test disturbance sample; the a-posteriori inflation uses a
separate validation sample.

\section{Results}
\label{sec:results}

\subsection{The selection effect is real and large}
The normalized back-off, marginally valid and adaptive in width, is the method a soft-sensor
interval would supply. Under optimization it realizes a violation of $0.429$ at the realistic
disturbance, rising to $0.488$ across the range, roughly five times the $0.10$ target
(Fig.~\ref{fig:violation}, Table~\ref{tab:regime}). The optimizer settles in the low-$D$ corner
near $R\approx2.83$, $D\approx33.7$, close to the deterministic optimum, precisely where the
adaptive margin under-covers the conditional quantile. The fixed margin over-violates more
modestly, between $0.11$ and $0.19$, and becomes infeasible beyond $\sigma_z=0.010$ as the spec
back-off it must add exceeds the achievable range. This is Theorem~\ref{thm:vacuous} in the data: marginal validity does not
transfer to the selected operating point.

\begin{figure}[t]
\centering
\includegraphics[width=\linewidth]{F3_violation_vs_sigma.png}
\caption{Realized constraint-violation rate versus disturbance magnitude $\sigma_z$ for the four
back-offs, with the nominal target $\alpha=0.10$ and the realistic $\sigma_z$ marked. Marginal
back-offs (fixed, normalized) violate well above target; CQR with a-posteriori calibration tracks
the oracle. Shaded band is the Monte-Carlo uncertainty.}
\label{fig:violation}
\end{figure}

\subsection{The oracle isolates the cause}
The oracle back-off, which knows the true conditional quantile, holds the realized violation at
$0.079$ at the realistic disturbance and near $\alpha$ throughout, at a setpoint that backs the
column off to $R=3.00$, $D=34.3$. Because the oracle satisfies the chance constraint by
construction while the marginal methods do not, the failure lies in the back-off, not in the
chance-constrained formulation \eqref{eq:chance}. This is the empirical counterpart of the
oracle's role in the theory.

\subsection{Conditional calibration restores safety at near-oracle profit}
The CQR back-off with a-posteriori calibration holds the realized violation between $0.063$ and
$0.105$ for $\sigma_z\le0.020$, within about $0.01$ of the target and of the oracle, backing the
setpoint from $D\approx34.3$ at $\sigma_z=0.006$ to $D\approx36.6$ at $\sigma_z=0.020$ as the
disturbance grows. Where the conditional quantile alone is slightly optimistic, the a-posteriori
step closes the gap: at $\sigma_z=0.006$ an inflation $\kappa^\star=1.53$ restores the target. This
is Theorem~\ref{thm:repair} in the data, with the CQR construction supplying conditional validity
and the a-posteriori step supplying the residual finite-sample correction.

\subsection{The inflation factor bounds the safe regime}
The method is not unconditionally safe, and the inflation factor is the signal
(Fig.~\ref{fig:kappa}). Because the calibration budget scales with the disturbance
($n_{\mathrm{cal}}=1500\max(1,\sigma_z/0.006)$), $\kappa^\star$ does not climb monotonically; it is
near unity where the conditional quantile is adequate, rises where it is locally optimistic
($\kappa^\star=1.53$ at $\sigma_z=0.006$ and $\kappa^\star=5.78$ at $\sigma_z=0.015$), and at
$\sigma_z=0.025$, more than four times the realistic disturbance, no finite inflation certifies the
target and the instance is declared infeasible. Its excursions and the terminal infeasibility both
mark where the conditional estimate can no longer certify safety, making data inadequacy a readable
property of the method rather than a silent failure, as Corollary~\ref{cor:kappa} predicts.

\begin{figure}[t]
\centering
\includegraphics[width=0.85\linewidth]{F7_kappa_vs_sigma.png}
\caption{A-posteriori inflation factor $\kappa^\star$ for the CQR method versus $\sigma_z$. The
margin is stretched (rising $\kappa^\star$) before the method becomes infeasible, making
$\kappa^\star$ a data-adequacy diagnostic.}
\label{fig:kappa}
\end{figure}

\subsection{Calibrated safety, not profit}
At the realistic disturbance every method lies within half a percent of the deterministic
optimum in profit (Table~\ref{tab:regime}); the constraint is barely active there, so there is
little profit to separate the methods. The value of the conditional method is therefore the
violation it removes, not any profit it adds: it converts a fivefold constraint breach into a
calibrated one at negligible economic cost.

\subsection{Twin validation and reproducibility}
The truth surface is validated by leave-one-feed-composition-out interpolation, fitting on all
$z\neq0.375$ and predicting the held-out $z=0.375$ slice, giving $R^2=0.917$ and a mean absolute
error of $0.0053$ in $x_B$, comfortably below the specification. Monte-Carlo estimation moves
individual violation rates by at most $\pm0.01$ to $0.02$ across platforms while leaving the regime
structure, the fivefold failure, and the data-starvation boundary intact.

\begin{table}[t]
\centering
\caption{Regime map over disturbance magnitude $\sigma_z$. For each method: operating profit
(USD\,h$^{-1}$), realized violation $\hat v$, a-posteriori inflation $\kappa^\star$, and setpoint
$(R^\star,D^\star)$. INF denotes an infeasible instance; the realistic disturbance is
$\sigma_z=0.006$.}
\label{tab:regime}
\begin{tabular}{@{}llrrrrr@{}}
\toprule
$\sigma_z$ & method & profit & $\hat v$ & $\kappa^\star$ & $R^\star$ & $D^\star$ \\
\midrule
0.006 & fixed       & 5381 & 0.167 & 1.00 & 3.00 & 34.1 \\
      & normalized  & 5394 & 0.449 & 1.00 & 2.83 & 33.7 \\
      & CQR         & 5376 & 0.101 & 1.53 & 3.00 & 34.3 \\
      & oracle      & 5374 & 0.079 & 1.00 & 3.00 & 34.3 \\
\midrule
0.010 & fixed       & 5356 & 0.112 & 1.00 & 2.93 & 34.8 \\
      & normalized  & 5396 & 0.488 & 1.00 & 2.93 & 33.7 \\
      & CQR         & 5335 & 0.066 & 1.00 & 2.60 & 35.3 \\
      & oracle      & 5351 & 0.106 & 1.00 & 2.80 & 34.9 \\
\midrule
0.015 & fixed       & ---  & INF   & ---  & ---  & ---  \\
      & normalized  & 5396 & 0.481 & 1.00 & 3.00 & 33.7 \\
      & CQR         & 5298 & 0.105 & 5.78 & 3.00 & 36.2 \\
      & oracle      & 5310 & 0.095 & 1.00 & 2.43 & 35.9 \\
\midrule
0.020 & fixed       & ---  & INF   & ---  & ---  & ---  \\
      & normalized  & 5394 & 0.483 & 1.00 & 3.00 & 33.7 \\
      & CQR         & 5282 & 0.094 & 1.00 & 2.37 & 36.6 \\
      & oracle      & 5288 & 0.087 & 1.00 & 2.20 & 36.5 \\
\midrule
0.025 & fixed       & ---  & INF   & ---  & ---  & ---  \\
      & normalized  & 5396 & 0.502 & 1.00 & 3.00 & 33.7 \\
      & CQR         & ---  & INF   & ---  & ---  & ---  \\
      & oracle      & 5271 & 0.105 & 1.00 & 2.23 & 36.9 \\
\bottomrule
\end{tabular}
\end{table}

\section{Discussion}
\label{sec:discussion}

\subsection{The design rule}
The result is narrow in its case study but general in mechanism. Nothing in
Theorem~\ref{thm:vacuous} is specific to distillation or to feed disturbance: it follows whenever
a conformal interval is placed inside an optimization. First, split and normalized conformal
back-offs guarantee coverage on average and are unsafe at the optimum. Second, the optimizer is a
selection operator, so conditional validity, not marginal validity, is the binding requirement.
The portable rule is to size the margin conditionally and to read the a-posteriori inflation as
the signal that the data can no longer certify the target.

\subsection{Limitations}
The plant is a simulation twin, rigorous but not a real column; the mechanism is structural and
should transfer, but the specific magnitudes are twin-specific. The guarantees of
Theorem~\ref{thm:repair} assume exchangeable calibration (Assumption~\ref{ass:exch}); a real feed
stream is autocorrelated and drifts, which breaks exchangeability. The weighted and
beyond-exchangeability conformal machinery \cite{tibshirani2019,barber2023,gibbs2021} lifts this
assumption and is the route to a plant guarantee under drift; developing it, together with a
closed-loop realization in which an online soft sensor supplies the constrained estimate and the
disturbance evolves in time, is the natural next step. Finally, the a-posteriori step of
part~(a) requires an independent calibration draw at the selected point; in deployment this is a
held-out window of operating data, and its size sets the residual gap of Theorem~\ref{thm:repair}.

\section{Conclusion}
\label{sec:conclusion}
Real-time optimization must subtract a back-off from a quality specification to protect a
constraint it cannot measure, and conformal prediction is the natural distribution-free way to
size it. We proved that the natural choice is unsafe under optimization: a marginally valid
back-off carries no coverage guarantee at the operating point the optimizer selects, a
constraint-space winner's curse. We then proved that conditional calibration restores a
finite-sample, distribution-free guarantee, exactly under a-posteriori recalibration and up to an
$O(\sqrt{d/n})$ term under uniform calibration, and that the required inflation factor is an
observable data-adequacy diagnostic. On a rigorous debutanizer twin the theory holds: marginal
back-offs realize fivefold violations, the conditional method tracks a plant-aware oracle at
near-oracle profit, and the inflation factor bounds the safe operating regime. The mechanism and
its remedy apply to any conformal interval used as a hard optimization constraint, which makes
conditional calibration a general design requirement for safe data-driven optimization.

\bibliographystyle{IEEEtran}
\bibliography{references}

\vspace{-0.8em}
\begin{IEEEbiographynophoto}{Bien Don Busico}
is a board-certified Chemical/Process Engineer affiliated with Mapua Malayan College Mindanao,
Davao City, Philippines. His research interests include distribution-free uncertainty
quantification, real-time optimization, and safe data-driven control for process systems.
\end{IEEEbiographynophoto}

\end{document}
